\chapter{Практика 6. Модель канала передачи}

\section{Задание}
Реализовать модель многолучевого канала с аддитивным белым гауссовым шумом.

\section{Теоретическое описание}
Модель многолучевого канала учитывает распространение сигнала по нескольким путям (лучам). Каждый луч характеризуется:
\begin{itemize}
    \item Задержкой $\tau_i$ относительно прямого луча
    \item Ослаблением $G_i$, зависящим от расстояния
\end{itemize}

Задержка рассчитывается по формуле:
\[
\tau_i = \frac{D_i - D_1}{c \cdot T_s}, \quad i = 1, 2, \ldots, N_B
\]

Ослабление сигнала в вакууме:
\[
G_i = \frac{c}{4\pi D_i f_0}
\]

После суммирования всех лучей добавляется аддитивный белый гауссов шум (АБГШ) с заданной спектральной плотностью мощности $N_0$.

\section{Реализация}

\begin{lstlisting}[language=Python, caption=Модель многолучевого канала]
class MultipathChannel:
    def __init__(self, fc=2.4e9, bandwidth=9e6, num_paths=3, n0_db=-100):
        self.fc = fc
        self.bandwidth = bandwidth
        self.num_paths = num_paths
        self.n0_db = n0_db
        self.c = 3e8
        self.Ts = 1.0 / bandwidth

    def propagate(self, tx_signal):
        distances = np.random.uniform(10, 500, self.num_paths)
        min_dist = np.min(distances)
        delays = []
        for d in distances:
            tau = (d - min_dist) / (self.c * self.Ts)
            delays.append(int(round(tau)))
        gains = []
        for d in distances:
            gain = self.c / (4 * np.pi * d * self.fc)
            gains.append(gain)
        L = len(tx_signal)
        rx_sum = np.zeros(L + max(delays), dtype=complex)
        for i in range(self.num_paths):
            shift = delays[i]
            gain = gains[i]
            shifted = np.zeros(L + shift, dtype=complex)
            shifted[shift:] = tx_signal * gain
            rx_sum[:len(shifted)] += shifted
        rx_signal = rx_sum[:L]
        n0_linear = 10 ** (self.n0_db / 10.0)
        noise_power = n0_linear * self.bandwidth
        noise_std = np.sqrt(noise_power / 2.0)
        noise = noise_std * (np.random.randn(L) + 1j * np.random.randn(L))
        rx_signal += noise
        return rx_signal
\end{lstlisting}

\section{Результат}
При количестве лучей $N_B = 3$ и мощности шума $N_0 = -100$ дБ сигнал на выходе канала имеет ту же длину, что и на входе, но содержит искажения от многолучевости и шума.

\section{Вывод}
Модель многолучевого канала с АБГШ успешно реализована. Она учитывает задержки и ослабления для каждого луча, а также аддитивный шум.